\documentclass[11pt,letterpaper]{article}
\usepackage[T1]{fontenc}
\usepackage{lmodern,microtype,amsmath,amssymb,amsthm,booktabs}
\usepackage[margin=0.88in,headheight=14pt]{geometry}
\usepackage{enumitem,fancyhdr,xcolor}
\usepackage[colorlinks=true,linkcolor=black,urlcolor=blue!45!black]{hyperref}
\hypersetup{pdftitle={Statistics and Predictive Modeling: Course Notes},pdfauthor={Prepared for Sean Rogers}}
\setlength{\parindent}{0pt}
\setlength{\parskip}{6pt}
\setlist{itemsep=3pt,topsep=5pt,leftmargin=*}
\pagestyle{fancy}
\fancyhf{}
\fancyhead[L]{\small Statistics and Predictive Modeling}
\fancyhead[R]{\small Lesson 1}
\fancyfoot[C]{\thepage}
\renewcommand{\headrulewidth}{0.3pt}
\newcommand{\E}{\mathbb E}
\newcommand{\Var}{\operatorname{Var}}
\newcommand{\Cov}{\operatorname{Cov}}
\newcommand{\SE}{\operatorname{SE}}
\newcommand{\SD}{\operatorname{SD}}
\newcommand{\MSE}{\operatorname{MSE}}
\newcommand{\RMSE}{\operatorname{RMSE}}
\newtheorem{theorem}{Theorem}[section]
\newtheorem{proposition}[theorem]{Proposition}
\theoremstyle{definition}
\newtheorem{definition}[theorem]{Definition}
\newtheorem{example}[theorem]{Worked example}
\theoremstyle{remark}
\newtheorem*{remark}{Remark}
\numberwithin{equation}{section}
\begin{document}
\thispagestyle{empty}
\begin{center}
{\Large\scshape Course Notes}\par\vspace{0.3in}
{\Huge Statistics and\\[6pt] Predictive Modeling}\par\vspace{0.2in}
{\large An exercise-first course for Sean Rogers}\par
\vspace{0.15in}
{\normalsize Cumulative edition \textbullet\ Lesson 1 \textbullet\ September 22, 2026}
\end{center}
\vspace{0.15in}
\hrule
\vspace{0.12in}
\textbf{Lesson 1. Sampling distributions and generalization.}
The central question is not merely how well a model performed on one dataset,
but what that performance tells us about its behavior on future data.
These notes connect sample means, uncertainty, paired comparisons, and hypothesis
tests to that question.

\textbf{How to use these notes.} Read the definitions and proofs, then cover the
worked solutions and reproduce their reasoning. The final sections contain
retrieval exercises and a separate answer key. The interactive lesson is designed
for approximately 30 minutes; this reference develops the details more fully.

\tableofcontents
\vfill
\textbf{Cumulative-document convention.} New completed lessons will be appended
as new numbered sections, preserving prior lessons and updating the contents.
This edition contains Lesson 1 only. The matching \texttt{.tex} source is the
editable master.
\newpage

\section{Lesson 1: Sampling distributions and generalization}
\subsection{The prediction problem and its statistical objects}
\begin{definition}[Population, sample, statistic, and parameter]
A \emph{population distribution} $P$ describes possible observations $Z=(X,Y)$,
where $X$ contains predictors and $Y$ is the outcome. A \emph{sample} is a
collection $Z_1,\ldots,Z_n$ drawn from that population. A \emph{statistic}
$T=T(Z_1,\ldots,Z_n)$ is a function of the observed data that does not depend
on unknown parameters. A \emph{parameter}, such as $\E_P[Y]$, is a fixed
feature of $P$ in the frequentist framework.
\end{definition}
\begin{definition}[Sampling distribution and standard error]
The \emph{sampling distribution} of $T$ is its probability distribution under
repeated sampling using the specified sampling procedure. Its \emph{standard
error} is $\SE(T)=\sqrt{\Var(T)}$. An estimated standard error replaces unknown
population quantities with estimates.
\end{definition}
The statistic is random \emph{before} the sample is observed; its realized value
is a number. A sampling distribution is not the distribution of individual
observations within one sample.

\textbf{Notation.} A fitted model is $\hat f$. Unless explicitly stated otherwise,
we hold $\hat f$ fixed and evaluate it using fresh observations independent of
its training data. Squared loss is $L=(Y-\hat f(X))^2$.

\begin{example}[Training performance versus generalization]
The original house-price model has training RMSE \$8,000 and test RMSE \$42,000.
A simpler model has training RMSE \$25,000 and test RMSE \$30,000.

\emph{Interpretation.} The first model's large gap is consistent with overfitting,
assuming a comparable train/test distribution and a valid evaluation. The simpler
model performs better on this test set. Its smaller gap does not prove underfitting:
underfitting concerns failure to capture useful structure, relative to attainable
performance, not merely the size of the gap.

\emph{What remains uncertain?} Results may depend on the test sample, distribution
shift, or prior model selection. Compare against a baseline, inspect learning
curves, and use validation appropriate to the intended deployment population.
If test results are used to choose a model, that set is functioning as validation
data; obtain a fresh final evaluation for the selected model.
\end{example}
\textbf{Terminology.} RMSE is an evaluation statistic. ``Test statistic'' usually
means a statistic used in a hypothesis test, not any statistic calculated on a
machine-learning test set.
\newpage

\subsection{Sample means: expectation, variance, and precision}
\begin{theorem}[Moments of an i.i.d. sample mean]
Let $X_1,\ldots,X_n$ be independent and identically distributed, with
$\E[X_i]=\mu$ and $\Var(X_i)=\sigma^2<\infty$. Then
\begin{equation}
\bar X=\frac1n\sum_{i=1}^nX_i,\qquad
\E[\bar X]=\mu,\qquad
\Var(\bar X)=\frac{\sigma^2}{n},\qquad
\SE(\bar X)=\frac{\sigma}{\sqrt n}.
\end{equation}
\end{theorem}
\begin{proof}
Linearity of expectation gives
\[
\E[\bar X]=\frac1n\sum_{i=1}^n\E[X_i]=\mu.
\]
Using $\Var(cX)=c^2\Var(X)$ and independence,
\[
\Var(\bar X)=\frac1{n^2}\Var\!\left(\sum_{i=1}^nX_i\right)
=\frac1{n^2}\sum_{i=1}^n\Var(X_i)=\frac{n\sigma^2}{n^2}.
\]
Taking the square root gives the standard error.
\end{proof}
Independence is not needed for the expectation calculation. Pairwise zero
covariances suffice for the displayed variance calculation. Normality is not
needed for either result.

\begin{example}[Mean house price]
Individual prices have mean \$400,000 and standard deviation \$150,000.
For an i.i.d. sample of 100 houses,
\[
\E[\bar X]=\$400{,}000,\qquad
\SE(\bar X)=\frac{\$150{,}000}{10}=\$15{,}000.
\]
Individual prices still have standard deviation \$150,000. It is the sample mean
whose standard deviation is \$15,000.
\end{example}
\begin{example}[Halving uncertainty]
To halve a standard error proportional to $n^{-1/2}$, solve
\[
\frac{\sigma}{\sqrt{n_{\mathrm{new}}}}
=\frac12\frac{\sigma}{\sqrt{n_{\mathrm{old}}}}
\quad\Longrightarrow\quad n_{\mathrm{new}}=4n_{\mathrm{old}}.
\]
This is exact for the sample mean under the stated assumptions. For more general
metrics, such as RMSE, it is typically an asymptotic rule requiring additional
regularity. It is not a guarantee for every metric or sampling design.
\end{example}
\textbf{When $\sigma$ is unknown.} With
$s^2=(n-1)^{-1}\sum_i(X_i-\bar X)^2$, the usual estimated standard error is
$s/\sqrt n$. Always distinguish a variance from its square root.
\newpage

\subsection{The law of large numbers and the central limit theorem}
These theorems answer different questions: \emph{where does an average go}, and
\emph{what do its rescaled fluctuations look like}?

\begin{theorem}[Weak law of large numbers; finite-variance version]
Under the sample-mean assumptions above, for every $\varepsilon>0$,
\[
\Pr(|\bar X_n-\mu|>\varepsilon)\longrightarrow0.
\]
Thus $\bar X_n\xrightarrow{p}\mu$.
\end{theorem}
\begin{proof}
Chebyshev's inequality gives
\[
\Pr(|\bar X_n-\mu|>\varepsilon)
\leq\frac{\Var(\bar X_n)}{\varepsilon^2}
=\frac{\sigma^2}{n\varepsilon^2}\longrightarrow0.
\]
\end{proof}

\begin{theorem}[Classical i.i.d. central limit theorem]
If the $X_i$ are i.i.d. with mean $\mu$ and variance $0<\sigma^2<\infty$, then
\begin{equation}
\frac{\sqrt n(\bar X_n-\mu)}{\sigma}\xrightarrow{d}N(0,1).
\end{equation}
\end{theorem}
For an adequate sample size, this motivates
$\bar X_n\approx N(\mu,\sigma^2/n)$. Here the second argument of $N$ denotes
\emph{variance}, not standard deviation. The CLT is stated without proof in this
lesson.

\textbf{What the CLT does not say.}
\begin{itemize}
\item Individual observations become normally distributed as $n$ increases.
\item Every sample of size 30 or 100 yields an accurate normal approximation.
\item Dependence, infinite variance, or selection bias can be ignored.
\end{itemize}
If individual observations are normal, their sample mean is exactly normal for
every $n$. Otherwise the approximation can be slow, especially for heavy tails.

\begin{example}[Approximate probability for an estimated MSE]
Suppose squared losses have mean 900 and variance 40,000. With $n=100$,
the sample MSE has mean 900 and standard error 20, so the CLT suggests
\[
\Pr(860\leq\widehat{\MSE}\leq940)
\approx\Pr\!\left(\frac{860-900}{20}\leq Z\leq\frac{940-900}{20}\right)
=\Pr(-2\leq Z\leq2)\approx0.9545.
\]
Both endpoints are two standard errors from the mean. This is a repeated-sample
probability with population mean specified, not yet an interval estimating an
unknown population mean.
\end{example}
\newpage

\subsection{Prediction risk, MSE, and RMSE}
\begin{definition}[Population risk and empirical test risk]
For a fixed model $\hat f$ and loss function $\ell$, the population prediction
risk is $R(\hat f)=\E_P[\ell(Y,\hat f(X))]$. Its empirical test estimate is
\[
\widehat R=\frac1n\sum_{i=1}^n\ell(Y_i,\hat f(X_i)).
\]
The expectation is over a new observation from the target population.
\end{definition}
For squared loss, set $L_i=(Y_i-\hat f(X_i))^2$. Then
\[
\widehat{\MSE}=\bar L,\qquad
\widehat{\RMSE}=\sqrt{\bar L},\qquad
\MSE=\E[L],\qquad \RMSE=\sqrt{\E[L]}.
\]
If test observations are i.i.d. and independent of training, conditional on the
fitted model the losses are i.i.d. With $\mu_L=\E[L]$ and
$\tau^2=\Var(L)<\infty$,
\begin{equation}
\E[\widehat{\MSE}]=\mu_L,\quad
\Var(\widehat{\MSE})=\frac{\tau^2}{n},\quad
\SE(\widehat{\MSE})=\frac{\tau}{\sqrt n}.
\end{equation}
\begin{example}[The full MSE calculation]
For $\mu_L=900$, $\tau^2=40{,}000$, and $n=100$,
\[
\E[\widehat{\MSE}]=900,\quad
\Var(\widehat{\MSE})=400,\quad
\SE(\widehat{\MSE})=20.
\]
The standard deviation of an individual squared loss is $\tau=200$; it is not
20. MSE and its standard error have units of outcome squared.
\end{example}
\textbf{Moment condition.} Finite variance of squared loss requires a finite
fourth moment of the prediction error: $\E[(Y-\hat f(X))^4]<\infty$.
Finite MSE alone does not ensure the usual CLT for estimated MSE.

\textbf{RMSE is nonlinear.} In general,
\[
\E[\sqrt{\bar L}]\neq\sqrt{\E[\bar L]}.
\]
Concavity of the square root implies a downward bias by Jensen's inequality.
For $\mu_L>0$, a first-order Taylor approximation (the \emph{delta method}) gives
\[
\SE(\widehat{\RMSE})\approx
\frac{\tau}{2\sqrt{\mu_L}\sqrt n}.
\]
In the example this is $20/(2\cdot30)=1/3$ outcome units. This is an approximation,
not an exact sample-mean formula for RMSE.

\textbf{Two sources of randomness.} Retraining on new training data changes
$\hat f$; drawing new test data changes its measured score. Our calculations
isolate the latter, conditional on the fitted model.
\newpage

\subsection{Paired comparisons: compare the models on the same cases}
For two fixed models evaluated on the same observation $i$, define
\[
L_{A,i}=(Y_i-\hat f_A(X_i))^2,\quad
L_{B,i}=(Y_i-\hat f_B(X_i))^2,\quad D_i=L_{A,i}-L_{B,i}.
\]
Let $\Delta=\E[D_i]=R(\hat f_A)-R(\hat f_B)$.
Negative $\Delta$ favors $A$; positive $\Delta$ favors $B$.
\begin{proposition}[Pairing preserves the estimand]
The average paired loss difference equals the difference of empirical MSEs:
\[
\bar D=\frac1n\sum_i(L_{A,i}-L_{B,i})
=\widehat{\MSE}_A-\widehat{\MSE}_B.
\]
Moreover,
\begin{equation}
\Var(D_i)=\Var(L_{A,i})+\Var(L_{B,i})-2\Cov(L_{A,i},L_{B,i}).
\end{equation}
\end{proposition}
\begin{proof}
The first identity is linearity of summation. For the second, expand the square
of $(L_A-\E L_A)-(L_B-\E L_B)$ and take expectations.
\end{proof}
Pairing does not demand that $A$ win on every case. It measures its \emph{average}
advantage while accounting for shared case difficulty. Positive within-case
covariance of the two losses reduces the variance of their difference.

\begin{example}[Paired standard error]
Suppose
\[
\Var(L_A)=400,\qquad\Var(L_B)=500,\qquad\Cov(L_A,L_B)=350.
\]
Then
\[
\Var(D_i)=400+500-2(350)=200.
\]
For 100 independent test cases,
\[
\Var(\bar D)=\frac{200}{100}=2,\qquad
\SE(\bar D)=\sqrt2\approx1.414.
\]
Ignoring the pairing would give $\sqrt{(400+500)/100}=3$, which is too large
for these paired data. Truly independent test sets of 100 cases per model
would instead have that standard error under these marginal variances.
\end{example}
\textbf{In practice.} Estimate the standard deviation of the observed $D_i$:
\[
s_D^2=\frac1{n-1}\sum_i(D_i-\bar D)^2,\qquad
\widehat{\SE}(\bar D)=\frac{s_D}{\sqrt n}.
\]
This is simpler than separately estimating all three variance/covariance terms.
The numerical variances in this example are a new hypothetical setup, separate
from the previous MSE example.
\newpage

\subsection{Confidence intervals and two-sided tests}
\begin{definition}[Frequentist confidence interval]
A confidence procedure $C(\text{data})$ with 95\% coverage contains the fixed
parameter in 95\% of repeated samples (approximately for an asymptotic procedure).
After observation, the realized interval either contains the parameter or does
not. Coverage is not a posterior probability for the parameter.
\end{definition}
If $\bar D$ is approximately normal with known standard error $s_0$, an approximate
95\% confidence interval for $\Delta$ is
\begin{equation}
\bar D\pm1.96s_0.
\end{equation}
To test $H_0:\Delta=0$ against $H_1:\Delta\neq0$, use
\[
Z_{\mathrm{obs}}=\frac{\bar D-0}{s_0},\qquad
p_{\mathrm{two}}=2\Phi(-|Z_{\mathrm{obs}}|),
\]
where $\Phi$ is the standard normal cumulative distribution function.
Reject at level $\alpha=0.05$ when $|Z_{\mathrm{obs}}|>1.96$.

\begin{example}[The observed advantage is $\bar D=-3$]
Using $s_0=\sqrt2$,
\[
Z_{\mathrm{obs}}=-\frac3{\sqrt2}\approx-2.1213,
\qquad p_{\mathrm{two}}\approx0.0339.
\]
The approximate 95\% interval is
\[
-3\pm1.96\sqrt2\approx(-5.772,-0.228).
\]
Zero is excluded, so the test rejects the equal-risk null at the 5\% level.
Your shortcut $-3\pm2\sqrt2\approx(-5.828,-0.172)$ gives the same decision;
under the normal approximation it has approximately 95.45\% coverage.
\end{example}
\textbf{Why the methods agree.} For any proposed value $\Delta_0$,
\[
\Delta_0\notin[\bar D-1.96s_0,\bar D+1.96s_0]
\quad\Longleftrightarrow\quad
\left|\frac{\bar D-\Delta_0}{s_0}\right|>1.96.
\]
Thus a two-sided confidence interval inverts the corresponding family of tests.

\textbf{What a $p$-value means.} It is a probability, under the null model, of a
test statistic at least as extreme in the specified tail(s). It is not the
probability that the null is true, nor the probability that the result was
``caused by chance.'' Statistical evidence is also distinct from practical
importance: a 3-unit squared-loss reduction needs outcome-scale context.

\textbf{Estimated standard errors.} Replacing $s_0$ with $s_D/\sqrt n$ gives a
paired $t$ statistic. Under i.i.d. normal $D_i$, it has an exact $t_{n-1}$ null
distribution. For nonnormal $D_i$, inference is generally approximate and can
be unreliable with small samples or heavy tails.
\newpage

\subsection{One-sided tests: direction must be chosen in advance}
If the prespecified question is whether model $A$ improves on model $B$, use
\[
H_0:\Delta\geq0\qquad\text{versus}\qquad H_1:\Delta<0.
\]
In the normal model with fixed known standard error, the largest lower-tail
rejection probability within this null occurs at its boundary $\Delta=0$.
Thus the lower-tail $p$-value is
\[
p_{\mathrm{one}}=\Phi(Z_{\mathrm{obs}}).
\]
At level 5\%, reject when $Z_{\mathrm{obs}}<-1.645$.

\begin{example}[One-sided inference for the same observed difference]
For $Z_{\mathrm{obs}}=-2.1213$,
\[
p_{\mathrm{one}}\approx0.0169.
\]
An equivalent 95\% one-sided confidence set is
\[
(-\infty,U],\qquad U=\bar D+1.645s_0
=-3+1.645\sqrt2\approx-0.674.
\]
Because $U<0$, the set excludes all nonnegative values of $\Delta$, supporting
the prespecified improvement claim.
\end{example}
\textbf{Why one-sided testing has more directional power.} It allocates all of
the 5\% rejection probability to the lower tail. A two-sided 5\% test allocates
2.5\% to each tail. A 95\% two-sided confidence interval is therefore not the
interval counterpart of a one-sided 5\% test.

\textbf{Why not choose the favorable direction afterward?} Under a continuous
symmetric null, selecting whichever tail matches the observed sign and rejecting
when its one-sided $p$-value is below 0.05 rejects in either 5\% tail. The actual
Type I error is 10\%, not 5\%.

\textbf{A surprising result still matters.} If the lower-tail test produces
$Z_{\mathrm{obs}}=+5$, it does not establish the chosen claim that $A$ is better.
It is strong evidence in the opposite direction, which should be reported.
Prespecifying the test does not mean ignoring harmful results.

\textbf{Decision guide.}
\begin{itemize}
\item Use a prespecified one-sided test when the inferential claim is directional,
such as a candidate improving on an incumbent.
\item Use a two-sided test when asking whether the models differ in either
direction, or when the direction was not prespecified.
\item Report effect size and uncertainty regardless of the test. Do not select
the test, direction, or model after looking at test outcomes and then report
an unadjusted confirmatory $p$-value.
\end{itemize}
\newpage

\subsection{Dependence, effective information, and validation design}
\begin{proposition}[Variance of an average with dependence]
For random variables with finite second moments,
\begin{equation}
\Var(\bar D)=\frac1{n^2}\left[
\sum_i\Var(D_i)+2\sum_{i<j}\Cov(D_i,D_j)\right].
\end{equation}
\end{proposition}
This follows by expanding the variance of a sum. Positive covariance between
different observations increases uncertainty relative to the independence formula.

\begin{example}[Ten neighborhoods, ten houses each]
Suppose the 100 test observations come from 10 independent neighborhoods.
If loss differences $D_i$ for houses within a neighborhood have positive
covariance, treating all 100 differences as independent underestimates
$\SE(\bar D)$: it drops positive terms from the formula above.
\end{example}
\textbf{Important precision correction.} The relevant dependence is between
the \emph{loss differences} $D_i$, not simply between raw prediction errors or
predicted house prices. Correlated residuals do not by themselves prove positive
covariance of squared-loss differences. Neighborhood effects might cancel
between models. The conclusion holds when the $D_i$ themselves have the assumed
positive within-neighborhood covariance.

\textbf{Preview: effective sample size.} Suppose every cluster has $m$ observations,
each $D_i$ has variance $\sigma_D^2$, within-cluster correlation is $\rho$, and
different clusters are independent. Counting correlated pairs yields
\[
\Var(\bar D)=\frac{\sigma_D^2}{n}\bigl[1+(m-1)\rho\bigr],
\qquad n_{\mathrm{eff}}=\frac{n}{1+(m-1)\rho}.
\]
For $n=100$, $m=10$, and $\rho=0.2$, the design effect is $2.8$,
$n_{\mathrm{eff}}\approx35.7$, and the true standard error is
$\sqrt{2.8}\approx1.67$ times the independence-based standard error.
This formula relies on the equal-size, equal-correlation assumptions.

\textbf{Two different covariance effects.}
\begin{itemize}
\item Positive covariance of $L_{A,i}$ and $L_{B,i}$ \emph{on the same case}
reduces the variance of their difference through the $-2\Cov$ term.
\item Positive covariance of $D_i$ and $D_j$ \emph{across different cases}
increases the variance of their average through the $+2\Cov$ terms.
\end{itemize}
\textbf{Choose validation to match deployment.} Holding out entire neighborhoods
assesses transfer to unseen neighborhoods. Predicting new houses in familiar
neighborhoods is a different target; a within-neighborhood split is not inherently
invalid for that goal. Dependence-aware uncertainty may still be required.
Cluster-level resampling or cluster-robust methods can help, but only 10 independent
clusters is a small sample and requires caution.
\newpage

\subsection{Retrieval practice and review plan}
Attempt these without referring to the worked examples. Explain your assumptions,
not just your arithmetic. Solutions begin on the next page.
\begin{enumerate}[label=\textbf{\arabic*.}]
\item Define a statistic, its sampling distribution, and its standard error.
What is held fixed when we evaluate one already-trained model repeatedly?
\item For i.i.d. $X_i$ with mean $\mu$ and variance $\sigma^2$, derive
$\E[\bar X]$ and $\Var(\bar X)$. Which step needs zero covariances?
\item You want one-third of your current standard error. How must $n$ change
under $n^{-1/2}$ scaling?
\item State the LLN and CLT in words. Which concerns consistency and which
describes scaled fluctuations?
\item If losses have mean 900 and variance 40,000, calculate the mean,
variance, and standard error of their average for $n=400$.
\item Explain why $\widehat{\MSE}$ is unbiased for population MSE but
$\widehat{\RMSE}$ is generally not unbiased for population RMSE.
\item Models are evaluated on the same cases. Their individual loss variances
are 400 and 500 and their within-case loss covariance is 350. Calculate
$\Var(D_i)$ and $\SE(\bar D)$ for $n=100$ independent cases.
\item For $\bar D=-3$ and known standard error $\sqrt2$, compute the $z$ statistic,
the two-sided 95\% interval, and the one-sided 95\% upper confidence bound.
Explain the associated hypotheses.
\item Why is it invalid to select a one-sided test's direction after seeing
the observed effect? What is the null rejection rate when either 5\% tail
is selected in this way under a symmetric continuous null?
\item Explain why positive within-case covariance helps a paired comparison
but positive across-case covariance can hurt precision. Identify the random
variables in each covariance.
\item A test set contains 10 independent neighborhoods with 10 houses each.
Under the equal-correlation model with $\rho=0.2$, find the effective sample
size. Would you use neighborhood-held-out validation for new neighborhoods
or necessarily for every prediction task?
\end{enumerate}

\textbf{Spaced-review plan.} On the next study day, retrieve Questions 2, 7, and 10.
About three days later, retrieve Questions 4, 6, and 8. After one week, mix all
topics without notes; after two weeks, revisit missed items and solve numerical
variants. These are study intervals, not additional scheduled notifications.

\textbf{Personal review priorities from this lesson.} Your conceptual reasoning
was strong: you recognized split uncertainty, understood paired comparisons,
and correctly used interval inversion. Revisit the factor of 2 in the covariance
identity, variance versus standard error, and one-sided versus two-sided thresholds.

\textbf{Next lesson.} Deeper LLN/CLT intuition, dependence and effective sample
size, and why more data can make a biased evaluation more precise without making
it more relevant.
\newpage

\subsection{Retrieval answer key and compact formula sheet}
\begin{enumerate}[label=\textbf{\arabic*.},itemsep=6pt]
\item A statistic is a function of sampled data; its sampling distribution is
its distribution across repeated samples; its standard error is its standard
deviation. For fixed-model evaluation, hold the fitted prediction rule fixed
and vary the fresh test observations.
\item $\E[\bar X]=\mu$ by linearity. In general,
$\Var(\bar X)=n^{-2}\sum_{i,j}\Cov(X_i,X_j)$.
With pairwise zero covariances and equal variances this is $\sigma^2/n$.
\item Multiply $n$ by 9.
\item LLN: the average approaches its population mean in probability under
the stated assumptions. CLT: its centered, $\sqrt n$-scaled fluctuations
converge in distribution to a normal law. LLN gives consistency; CLT supports
approximate uncertainty calculations.
\item Mean 900; variance $40{,}000/400=100$; standard error 10.
\item MSE is an average of losses, so expectation is preserved. The square
root is nonlinear and concave, so $\E[\sqrt{\bar L}]\leq\sqrt{\E[\bar L]}$.
\item $\Var(D_i)=400+500-700=200$; $\SE(\bar D)=\sqrt{200/100}=\sqrt2$.
\item $z=-2.1213$; two-sided interval $(-5.772,-0.228)$; one-sided upper
bound $-0.674$. The two-sided test asks whether $\Delta\neq0$; the lower-tail
test asks whether $\Delta<0$ against the composite null $\Delta\geq0$.
\item Choosing the observed direction creates a rejection region comprising
both 5\% tails: total Type I error 10\%, not the claimed 5\%.
\item $\Cov(L_{A,i},L_{B,i})$ enters $\Var(D_i)$ with coefficient $-2$.
$\Cov(D_i,D_j)$ enters the variance of the average with a positive coefficient.
They are different covariances, so there is no contradiction.
\item $n_{\mathrm{eff}}=100/[1+9(0.2)]\approx35.7$. Group-held-out validation
targets new neighborhoods; validation for familiar neighborhoods should match
that different deployment setting.
\end{enumerate}
\vspace{3pt}\hrule\vspace{5pt}
\textbf{Formula sheet (under the stated assumptions)}
\begin{align*}
\SE(\bar X)&=\sigma/\sqrt n,
&\widehat{\SE}(\bar X)&=s/\sqrt n,\\
\Var(A-B)&=\Var(A)+\Var(B)-2\Cov(A,B),\\
Z&=(\widehat\theta-\theta_0)/\SE(\widehat\theta),
&\mathrm{CI}_{95\%}&:\ \widehat\theta\pm1.96\SE(\widehat\theta),\\
p_{\mathrm{two}}&=2\Phi(-|Z|),
&p_{\mathrm{lower}}&=\Phi(Z).
\end{align*}
\textbf{Always check:} target population; fixed versus refitted model; independence
or clustering; finite moments; untouched evaluation data; prespecified test
direction; practical effect size.
\end{document}